\scanpage{24}
\deleted{
\[
(5)\qquad\Delta_A(t)=\Delta_{|A|}(t)=\Delta_{\varphi_{|A|}}(t)=\exp\int_0^t\log\varphi_{|A|}(s)\,ds
\]
\[
(6)\qquad\Delta(A)=\deleted{\int_0^{T(1)}}\Delta_A(\omega)\qquad(\text{si }\omega=T(1)<+\infty)
\]
Par suite, $\Delta_A(t)$ est une fonction \deleted{\unsure{finie}} continue \deleted{\unsure{type|exp}} de $t$ à valeurs $0\leq\Delta_A(t)<\infty$. \deleted{$\Delta(A)$ est fini si}

Si $T(1)<\infty$ (i.e. si $T$ est \unsure{une trace finie}), \unsure{alors|sinon} si la trace est finie, i.e. $1\leq\Delta(A)<+\infty$.
}

\margin{Définition de \deleted{$\varphi_{1+f}$} $\varphi_f$, $f\in\bar b$ $(\bar b=C+\underline b)$.}
\margin{Définition de $W(A)$. les ``Aux.'' C'est \deleted{\unsure{défini}} \unsure{par} $\unsure{\varphi}$ \unsure{en} $\unsure{\bar b}$.}

Supposons que $\underline c$ soit une sous-algèbre commutative auto-adjointe maximale dans $\underline b$, $\underline c=C_0(M)$, $\unsure{\mu|\rho}$ la mesure sur $M$ définie par le corollaire 1. Si $t\in M$ est un pt de $M$ de mesure non nulle, on sait (\unsure{par la condition de maximalité} de $\underline c$) que \inserted{$E=U_{\varphi_{\{t\}}}$} \deleted{$E_\varphi$} est un projecteur minimal de \deleted{$A\ \underline b$} $A$. Réciproquement soit $E$ est un projecteur minimal dans \deleted{$\unsure{\underline b}$}\inserted{$A$},
\deleted{
$\underline c$ \inserted{$=C_0(M)$} \unsure{sous-algèbre} \unsure{auto-adjointe} maximale de $\underline b$, contenant $E$, alors \deleted{il} s'\unsure{identifie} à la fonction.

$E$ est la fonction caractéristique d'un \deleted{\illegible{1}} \unsure{ensemble à la fois} \deleted{vide et fermé} \unsure{ouvert et fermé} de $M$, \unsure{donc forcément réduit} à un point \deleted{\unsure{isolé}} \unsure{de} \unsure{mesure} $>0$. On sait que $E$ Tout \deleted{\illegible{1}} caractéristique \unsure{d'un}

un projecteur minimal de $\underline b$, i.e. si \deleted{\illegible{2}} la mesure \unsure{sur} un \unsure{espace} $M$. Cela \deleted{\illegible{4}} tout $\unsure{C_0(M)}$ isolé \unsure{dans} notre algèbre $A$.
}
\inserted{(Si $E$ est un projecteur \inserted{$>0$} \deleted{\unsure{de}} minimal, \deleted{\unsure{dans}} on voit \unsure{facilement} qu'il \unsure{existe|est} un pt}
\mspage{25}
